\section{Conclusion}
\label{sec:conclusion}

We presented \method{}, a framework built on a simple but consequential claim:
traffic accident anticipation should not be viewed only as frame-wise
classification with optional explanation, but can also be studied through a
\emph{concept-conditioned sequential-decision lens}. In the current paper, the
strongest empirical support is for the WHY layer and for the value of a
concept-conditioned timing formulation; the policy-level WHEN story remains
only partially validated.

The central architectural consequence is to place concept activations directly
inside the policy state, $\mathbf{s}_t=[\mathbf{h}_t\|\mathbf{c}_t]$.
This turns semantic evidence into part of the decision pathway and makes the
timing mechanism structurally concept-conditioned, even though the current
headline metrics are still computed from the classifier trajectory rather than a
fully validated standalone actor policy.

A second contribution of this work is that the semantic interface is no longer
treated as a fixed handcrafted assumption. We introduced a reproducible
multimodal risk concept discovery and adjudication pipeline that mines driving
frames, performs risk-first refinement, and exports trainable ontologies for
the concept bottleneck. This matters because a CBM is only as useful as the
concept vocabulary it exposes. By explicitly constructing and comparing
historical, manual, and discovered concept sets, the paper elevates ontology
construction to the same scientific status as architecture design.

Empirically, \method{} achieves strong results on both benchmarks.
On A3D, it improves mTTA by +14.8\% over CRASH while maintaining
competitive AP---consistent with concept-conditioned timing being most
valuable when hazards develop progressively before impact. Matched
classifier-vs-actor comparisons further show that actor-based evaluation is
already fairly close to the classifier line on a representative A3D
checkpoint.
On DAD, \method{} delivers competitive performance within the
interpretable comparison tier, with multi-seed curriculum runs confirming
that prediction quality is reproducible at synchronized checkpoints even though
CRS robustness is weaker than predictive robustness. The same matched
comparison also shows that actor-based triggering on DAD remains substantially
weaker than classifier-based anticipation, which is why the paper continues to
report classifier-derived headline timing metrics.
Critically, these results are achieved with only 30M parameters---showing
that interpretability, parameter efficiency, and performance can be
mutually reinforcing rather than competing. We further show that compact
discovered concept sets yield non-trivial DAD/A3D operating-point changes under
matched launchers, supporting the claim that discovered concepts are trainable
primitives rather than appendix-only descriptive artifacts.
Finally, our current intervention package supports a more careful but still
important claim: the semantic interface is not merely descriptive. In archived
DAD analyses, positive concept amplification moves alerts earlier in a
measurable subset of valid edited cases (14\% earlier, 84\% unchanged, 2\%
later under the hybrid boost setting), although a stronger fully policy-level
intervention study remains future work.

\paragraph{Limitations and future work.}
First, concept labels are derived from CLIP pseudo-annotations:
high-quality human-labelled concept supervision would further strengthen
the WHY layer and reduce vocabulary noise. More importantly, exposed concept
channels should not automatically be read as concept correctness; the current
paper demonstrates architectural transparency and limited semantic auditing,
not a complete human-verified concept benchmark.
Second, the current evaluation uses offline VGG16 features;
end-to-end training with a modern visual backbone (e.g., ResNet-50 or ViT)
is a natural and important extension.
Third, the CAAC policy optimises a scalar time-remaining reward;
a multi-objective policy that jointly balances AP, mTTA, and
false-alarm rate could better expose deployment trade-offs and
support jurisdiction-specific calibration.
Fourth, while we now construct concept ontologies explicitly, broader
cross-dataset ontology transfer and stronger multi-seed ontology comparisons
remain open.
Fifth, a complete policy-level intervention analysis---re-running
the full Actor--Critic from edited concept states---remains an
open validation target that would further strengthen the causal
interpretability claim.
Sixth, the strongest ontology claim in the current paper is about reproducible
construction together with matched DAD/A3D operating-point comparisons; a
broader shared-recipe multi-seed ontology study would be the next major
empirical strengthening step.
Seventh, CRS should currently be treated as a tentative audit signal rather
than as a mature human-override interface, because clean-seed audits still show
frequent degeneracy.

\paragraph{Broader impact.}
\method{} is designed for safety-critical deployment settings such as
Advanced Driver Assistance Systems (ADAS) and autonomous driving research.
Its main practical value is to expose a more auditable intermediate semantic
state than a purely black-box anticipation model, which could in principle help
engineers inspect failure modes and reason about warning behavior. At the same
time, the current paper does \emph{not} establish that exposed concept channels
or CRS rankings are already reliable enough for direct human override in a real
deployment pipeline. As with all accident anticipation systems, false positives
(unnecessary alerts) must be balanced against false negatives (missed
accidents), and any real deployment would require substantially stronger
robustness, calibration, and human-factors validation than what is shown here.
